\begin{table}[ht]
\centering
\caption{The socioeconomic gradient in the compulsory-voting effect,
conditioning on geography (age-18 cutoff).}
\label{tab:gradient_geography}
\resizebox{\ifdim\width>\linewidth\linewidth\else\width\fi}{!}{%
\begin{tabular}{lrlrlrlrlrr}
\toprule
& \multicolumn{4}{c}{PCA index} & \multicolumn{4}{c}{NBI} & & \\
\cmidrule(lr){2-5}\cmidrule(lr){6-9}
Specification & \multicolumn{2}{c}{Men} & \multicolumn{2}{c}{Women}
              & \multicolumn{2}{c}{Men} & \multicolumn{2}{c}{Women}
              & $N$ & Circuits \\
\midrule
Unconditional & $-12.65^{**}$ & (1.52) & $-17.17^{**}$ & (1.51) & $-8.65^{**}$ & (1.39) & $-11.92^{**}$ & (1.37) & 214,710 & 4,914 \\
+ province fixed effects & $-10.57^{**}$ & (1.57) & $-15.03^{**}$ & (1.61) & $-7.61^{**}$ & (1.45) & $-11.34^{**}$ & (1.42) & 214,710 & 4,914 \\
Unconditional, density subsample & $-12.65^{**}$ & (1.52) & $-17.18^{**}$ & (1.51) & $-8.59^{**}$ & (1.39) & $-11.88^{**}$ & (1.37) & 214,391 & 4,866 \\
+ log voter density & $-12.88^{**}$ & (1.58) & $-17.41^{**}$ & (1.60) & $-7.94^{**}$ & (1.40) & $-11.25^{**}$ & (1.39) & 214,391 & 4,866 \\
Urban half of voters & $-9.83^{**}$ & (2.18) & $-13.74^{**}$ & (2.34) & $-5.12^{*}$ & (2.22) & $-9.48^{**}$ & (2.14) & 106,913 & 1,189 \\
Rural half of voters & $-14.77^{**}$ & (2.08) & $-19.73^{**}$ & (2.08) & $-10.44^{**}$ & (1.81) & $-12.73^{**}$ & (1.86) & 107,478 & 3,677 \\
\bottomrule
\end{tabular}%
}
\par\vspace{0.6em}
{\footnotesize\begin{minipage}{\linewidth}\emph{Note:} Each cell is
$\Delta\tau$, the change in the estimated cutoff effect from the 10th to the
90th percentile of the deprivation measure, in percentage points, with its
standard error in parentheses. Every row is the same continuous model as
Table~\ref{tab:triple_cells_T18}, fit in the same MSE-optimal window
($\approx$66 days), with standard errors clustered by electoral circuit (CRV1).
Each conditioning variable is interacted with $(1, x, T, Tx)$, so the
discontinuity itself is free to differ across provinces and across levels of
density. Urban and rural halves split at the median log voter density, so each
holds half the voters in the window; the urban half holds them in far fewer
circuits. Density is unavailable for $0.15\%$ of the window, so the last three
rows are read against the subsample row above them rather than against the
first. The first row reproduces the estimate reported in
Section~\ref{sec:res_ses}. $^{*}p<.05$, $^{**}p<.01$.
\end{minipage}}
\end{table}
